% =============================================================================
% 05_EVALUATION_METHODS — Analytical Synthesis of Evaluation Approaches
% =============================================================================
% Page target: ~2 pages
% Structure: Retrospective evaluation → Prospective evaluation → Continuous monitoring
% =============================================================================

\section{Evaluation Methods: From Retrospective Validation to Lifecycle Monitoring}
\label{sec:evaluation_methods}
\addcontentsline{toc}{section}{Evaluation Methods}

The evaluation of health AI systems occupies a spectrum ranging from
single-site retrospective validation — the current norm — to prospective,
longitudinal, ecologically grounded evaluation designs, and ultimately to
continuous post-market monitoring frameworks. This section traces that arc
analytically, showing how the field's evaluation practice has evolved, where
it remains structurally inadequate, and what emerging methodological responses
suggest about the requirements for a genuinely fit-for-purpose evaluation
ecosystem.

\subsection{Retrospective Evaluation: The Dominant Paradigm and Its Limits}
\label{sec:retrospective}

The prevailing evaluation paradigm in health AI relies on retrospective
analysis of curated datasets, split into training, validation, and test sets
within a single institution or health system. Performance is reported as
aggregate metrics — primarily Area Under the Receiver Operating Characteristic
Curve (AUC), accuracy, sensitivity, and specificity — computed on held-out
test data from the same population as the training data.

This approach has produced a structurally misleading evidence base.
\citet{abulibdeh2025illusion} document the consequences with particular
force: the majority of FDA-authorized AI devices were approved on the
basis of single-site, premarket studies with no mandatory external validation
and no continuous post-market monitoring. The result is what they term an
\textbf{illusion of safety}: statistically precise performance estimates that
conceal ecologically invalid assumptions about generalizability.

The MSI colorectal cancer prediction study exemplifies the pattern
quantitatively. Models reported AUCs of 0.882 in the training cohort,
declining to 0.768 in the internal validation cohort, 0.803 in external
validation cohort 1, and 0.751 in external validation cohort 2. This
progressive performance degradation across increasingly distant populations
is not an aberration — it is the expected pattern when evaluation conditions
do not reflect the complexity of real clinical deployment.

Three structural failures underpin this inadequacy. First, single-site
retrospective evaluation captures neither the temporal dynamics of clinical
decision-making nor the multi-professional workflows in which AI systems
are embedded \citep{turchi2025ecological}. A single-point-in-time accuracy
metric cannot represent the iterative, contextual, and socially situated
nature of clinical reasoning. Second, standard evaluation metrics collapse
qualitatively different clinical harms into a single number. A false-negative
in cancer diagnosis carries a fundamentally different clinical weight than
a false-positive, and the aggregated AUC obscures this asymmetry
\citep{turchi2025ecological}. Third, retrospective evaluation is
inherently incapable of detecting algorithmic drift: the degradation of
model performance over time as input data distributions shift due to
changes in patient populations, clinical practice, or the AI system's
own continuous learning updates.

\subsection{Prospective Longitudinal Evaluation: The Emerging Alternative}
\label{sec:prospective}

The methodological response to retrospective inadequacy has produced a
growing body of prospective, multicenter evaluation designs. The TRAC-Pain
protocol \citep{bailes2026trac} represents the most fully articulated
example: a 12-week prospective observational cohort study combining
continuous wearable sensor data — physiological, sleep, and physical
activity metrics — with ecological momentary assessments (EMAs)
administered multiple times daily. The design explicitly targets the
temporal dynamics and contextual complexity that retrospective evaluation
ignores, capturing the longitudinal pain experience rather than a
single point-of-care assessment.

This design represents a significant methodological advance along three
dimensions simultaneously. First, the continuous wearable data stream
provides objective, real-time performance signals rather than post-hoc
accuracy metrics. Second, the EMA component captures patient-reported
outcomes and contextual factors — pain interference, mood, fatigue —
that are invisible to model performance metrics but central to clinical
relevance. Third, the 12-week duration enables detection of performance
evolution over time, addressing the drift detection problem that
retrospective designs structurally cannot solve.

However, prospective longitudinal evaluation of this kind remains the
exception. The practical barriers are substantial: such designs are
costly, time-consuming, require sustained institutional commitment,
and are poorly aligned with the publication timelines and funding cycles
that incentivize academic health AI research. The result is a field in
which the methodological ideal is well understood but rarely implemented.

\subsection{Continuous Monitoring and Drift Detection}
\label{sec:continuous_monitoring}

Between the extremes of single-shot retrospective evaluation and
resource-intensive prospective cohort studies lies the operational
requirement for continuous post-market monitoring of deployed AI systems.
This is the dimension of evaluation most urgently under-developed in both
practice and literature.

\citet{warraich2024fda} identify continuous learning AI systems as requiring
special regulatory attention precisely because the traditional pre-market
validation paradigm cannot guarantee post-deployment performance. Yet mandatory
post-market performance monitoring with standardized drift detection protocols
is conspicuously absent from most regulatory frameworks reviewed. The FDA's
adverse event reporting mechanism captures safety incidents after they occur
but provides no proactive monitoring capability.

The gap between this need and existing practice is addressed most directly
by \citet{palama2026auditing}, whose Layer 3 (performance, safety, and drift
monitoring) specifies the operational architecture for continuous monitoring:
statistical process control methodologies applied to real-time model outputs,
with alert thresholds and predefined escalation pathways tied to institutional
decision authorities. The framework specifies that performance below a defined
operational threshold triggers an automatic governance notification, cascading
through defined stages from technical retraining to clinical review to market
suspension.

This operationalization is methodologically significant for two reasons.
First, it reframes evaluation as an organizational process rather than a
technical measurement event: the question is not only ``how accurate is
the model?'' but ``who is responsible for knowing how accurate the model
is at any given moment, and what must they do when it degrades?'' Second,
it shifts the unit of evaluation from model performance in isolation to
the sociotechnical system — model, clinicians, governance structure, and
institutional infrastructure — within which performance is sustained.

The ecological validity gap in evaluation methodology thus has a direct
governance consequence: a model validated under ecologically invalid
conditions generates misleading confidence estimates, which are then used
by governance bodies to justify deployment decisions that the evaluation
evidence does not support. Closing the ecological validity gap is not
merely a methodological improvement — it is a precondition for governance
that is honestly matched to the evidence.